\section{Some Useful Spaces}
%
\begin{definition}	\label{def:SpaceL2}
	Let \glssymbol{spaceL2} denote \gls{spaceL2}. Then
	\begin{equation}
		\glssymbol{spaceL2} = \left\{ \varphi \left( x \right) \bigg \vert \int_0^1 \left|\varphi \left( x \right)\right|^2 \dd x < \infty \right\}.
	\end{equation}
\end{definition}
%
\begin{definition}	\label{def:Space_H1}
	Let \glssymbol{spaceH1} denote \gls{spaceH1}. Then
	\begin{equation}
		\glssymbol{spaceH1} = \left\{ \varphi \left( x \right) \Big \vert \varphi \left( x \right) \in \glssymbol{spaceL2}, \;
			\frac{d \varphi }{dx} \in \glssymbol{spaceL2} \right\}.
	\end{equation}
\end{definition}
%
\begin{definition}	\label{def:Space_H10}
	Let \glssymbol{spaceH10} denote a \gls{spaceH10}. Then
	\begin{equation}
		\glssymbol{spaceH10} = \left\{ \varphi \left( x \right) \Big \vert \varphi \left( x \right) \in \glssymbol{spaceH1}, \;
			\varphi \left( 0 \right) = \varphi \left( 1 \right) = 0 \right\}.
	\end{equation}
\end{definition}
%
\subsection{Inner Products and Norms}
%
\begin{definition}	\label{def:L2_Inner_product}
	Let \glssymbol{L2InnerProduct} denote \gls{L2InnerProduct}. Then
	\begin{subequations}
		\begin{equation}	\label{eqn:L2_Inner_product}
			\glssymbol{L2InnerProduct} = \int_0^1 f \left( x \right) g \left( x \right) \dd x,
		\end{equation}
		and the natural norm
		\begin{equation}	\label{eqn:L2_Norm}
			\glssymbol{L2Normf} = \sqrt{ \left \langle f, f \right \rangle }.
		\end{equation}
	\end{subequations}	
\end{definition}
\begin{lemma}	\label{lem:Cauchy-Schwarz}
	Norms $\left\| f \right\|_{ L^2 }$ and $\left\| g \right\|_{ L^2 }$ obey Cauchy-Schwarz Inequality, i.e.,
	\begin{equation}	\label{eqn:Cauchy_Schwarz}
		\vert \glssymbol{L2InnerProduct} \vert \leq \glssymbol{L2Normf} \glsentrysymbol{L2Normg},
			\qquad \forall f, g \in \glssymbol{spaceL2}.
	\end{equation}
\end{lemma}
We also define a separate norm for functions in \gls{spaceH1}.
\begin{definition}	\label{def:H1_Norm}
	Let $\left\| f^\prime \right\|_{ H^1 }$ denote \gls{H1Normf}. Then
	\begin{equation}	\label{eqn:H1Norm}
		\glssymbol{H1Normf} = \sqrt{ \glssymbol{L2Normf}^2 + \big\| f^\prime \big\|_{ L^2 }^2 }.
	\end{equation}
\end{definition}
Note that while $\left\| f^\prime \right\|_{ L^2 }$ only measures the function magnitude, $\left\| f^\prime \right\|_{ H^1 }$ is a
measure of both function itself and its first derivative. We may also define a measure for \glsentryname{L2Norm} of the first derivative:
\begin{definition}	\label{def:H1_Semi_Norm}
	Let $\vert f \vert_{ H^1 }$ denote \gls{H1SemiNormf}. Then
	\begin{equation}	\label{eqn:H1_Semi_Norm}
		\glssymbol{H1SemiNormf} = \left\| f^\prime \right\|_{ L^2 }.
	\end{equation}
\end{definition}
To measure magnitude of the second derivative we define \glssymbol{H2SemiNormf}.
\begin{definition}	\label{def:H2_Semi_Norm}
	Let $\vert f \vert_{ H^2 }$ denote \gls{H2SemiNormf}. Then
	\begin{equation}	\label{eqn:H2_Semi_Norm}
		\glssymbol{H2SemiNormf} = \big\| f^{ \prime \prime }\big\|_{ L^2 } .
	\end{equation}
\end{definition}
\begin{remark}
	A \emph{true} norm is only zero when its input is zero. Norms $\vert f \vert_{ H^1 }$ and $\vert f \vert_{ H^2 }$ can give zero
	for non-zero inputs and are therefore called semi-norms.\footnote{For example for a linear polynomial function $f = ax + b$, we
	have $\vert f \vert_{ H^2 } = 0$.} In all other aspects they behave (mathematically) as norms.
\end{remark}
Finally we define \gls{aForm}
\begin{definition}	\label{def:a_Form}
	Let \glssymbol{aForm} denote \gls{aForm}. Then
	\begin{equation}	\label{eqn:a_Norm}
		\glssymbol{aForm} = \int_0^1 f^{\prime} \left(x \right) g^{\prime} \left(x \right) \dd x.
	\end{equation}
\end{definition}
\begin{remark}
	Note that \glssymbol{aForm} is a kind of inner product. It gives the $L^2$-Inner Product of its arguments' derivatives.
	Thus, for $f = g$, we have
	\begin{equation}
		a \left( f, f \right) = \int_0^1 \Bigl[ f^{\prime} \left(x \right) \Bigr]^2 \dd x = \glssymbol{H1SemiNormf}^2.
	\end{equation}
\end{remark}
%